\section{Rendered UX and Browser-Step QA}

The browser is a real runtime, not a screenshot afterthought. Agents can create UI that compiles, passes component tests, and still fails through overlapping text, broken responsive layout, invisible controls, stale selectors, missing focus states, blank canvases, cramped buttons, clipped labels, sticky footer collisions, or layout shift. Playwright's official guidance emphasizes user-visible behavior, resilient locators, isolation, and web-first assertions, which fit agent repair better than brittle sleeps and CSS selectors \cite{playwrightBestPractices2026,playwrightWritingTests2026}.

Rendered UX proof should be a contract, not a taste debate. The standard treats rendered UX as a layered proof system: design-system constraints, generated contracts, deterministic Storybook states, Playwright screenshots, DOM geometry rules, accessibility rules, visual regression, AI/CV triage, and human feedback labels. \texttt{@jankurai/ux-qa} is a reference implementation of the UX evidence contract, not the only conformant implementation.

\begin{lstlisting}[style=codebox]
{
  "route": "/checkout",
  "viewport": "390x844",
  "browser": "chromium-123",
  "screenshot_hash": "sha256:...",
  "aria_snapshot_hash": "sha256:...",
  "rule_ids": ["HLT-013", "UX-OVERLAP"],
  "selector": "[data-testid=submit-order]",
  "severity": "high",
  "decision": "block"
}
\end{lstlisting}

\begin{table}[t]
\caption{Required rendered-UX evidence fields.}
\label{tab:ux-evidence-fields}
\centering
\scriptsize
\begin{tabularx}{\columnwidth}{L{0.34\columnwidth} Y}
\toprule
\JKTableHeader
\textbf{Field group} & \textbf{Examples} \\
\midrule
Surface identity & Route, story, selector, viewport, browser, locale, timezone. \\
Artifact identity & Screenshot, crop, ARIA snapshot, trace, video, geometry report, hashes. \\
Rule evidence & Rule IDs, severity, deterministic detector result, baseline reference. \\
Privacy state & \texttt{clean}, \texttt{redacted}, \texttt{restricted}, or \texttt{blocked}. \\
Decision gate & \texttt{block} for objective failures; \texttt{review} for semantic visual judgment. \\
\bottomrule
\end{tabularx}
\end{table}

\begin{table*}[t]
\caption{Rendered UX QA detectors for a conformant product surface.}
\label{tab:rendered-ux-qa}
\centering
\small
\begin{tabularx}{\textwidth}{L{0.19\textwidth} Y Y}
\toprule
\JKTableHeader
\textbf{Detector} & \textbf{Machine signal} & \textbf{Merge behavior} \\
\midrule
Design-system constraints & Component imports, token use, no raw one-off spacing/color/z-index, Code Connect or equivalent mapping. & Block unapproved design-language drift. \\
Storybook state coverage & Normal, loading, empty, error, long text, locale, theme, mobile, role, and interaction states. & Require stories for reusable components and key screens. \\
Playwright visual proof & Component and flow screenshots under pinned browser, font, viewport, timezone, and locale settings. & Require review for protected baselines. \\
DOM geometry rules & Bounding boxes, computed styles, scroll sizes, target spacing, overlap, clipping, wrapping, nested scrollbars, and fixed/sticky collisions. & Block objective layout failures before human QA. \\
Accessibility structure & axe/WCAG checks, ARIA snapshots, keyboard paths, focus visibility, form labels, and target size. & Block critical known violations; route ambiguous cases to expert review. \\
Layout stability & CLS and late-shift evidence from Lighthouse or Web Vitals. & Block severe movement on product-critical pages. \\
AI/CV triage & Semantic review of ambiguous visual deltas and screenshot crops. & Warn or route to humans; do not replace deterministic gates. \\
Human feedback loop & Labels for true defect, intentional change, false positive, acceptable variance, missing rule, or design-system bug. & Improve thresholds, baselines, stories, and future rules. \\
\bottomrule
\end{tabularx}
\end{table*}

Some criteria are deterministic enough to be Jankurai evidence: overlap, clipping, horizontal overflow, target-size failures, missing focus visibility, severe CLS, blank critical surface, missing protected-route screenshot, clipped required text, sticky obstruction, and unapproved token violations. Ambiguous semantic visual changes, brand judgment, and acceptable variance route to review, not automatic pass. WCAG 2.2 target-size guidance defines a 24 by 24 CSS-pixel minimum or sufficient spacing exceptions for pointer targets \cite{wcagTargetSizeMinimum2026}. Core Web Vitals guidance keeps CLS of 0.1 or less as the good visual-stability threshold \cite{webDevCls2025}.

Open-source tools already solve real pieces. Storybook can turn every story into a visual baseline \cite{storybookVisualTests2026}. Playwright can compare screenshots and ARIA snapshots \cite{playwrightVisualComparisons2026,playwrightAriaSnapshots2026}. Argos and BackstopJS provide visual review and screenshot-diff workflows \cite{argosVisualTesting2026,backstopJs2026}. Playwright and Storybook can run axe-based accessibility checks, while their docs warn that automation is only a first-line signal \cite{playwrightAccessibility2026,storybookAccessibility2026}. Jankurai's role is to bind those outputs to route, viewport, browser, hashes, rule IDs, selector, severity, owner, and witness decision.

Pixel QA is first-class but not exhaustive on every pull request. Critical flows, high-risk UI surfaces, visual diff baselines, payments, auth, admin actions, onboarding, and canvas/3D surfaces need browser-step proof in the PR lane. Full viewport and device matrices belong in nightly or release lanes unless the changed path directly touches those surfaces.
